\documentclass[12pt]{article}
\usepackage[utf8]{inputenc}
\usepackage[T1]{fontenc}
\usepackage[margin=1in]{geometry}
\usepackage{setspace}
\usepackage{xcolor}
\usepackage[colorlinks=true, urlcolor=blue!60!black, linkcolor=blue!60!black]{hyperref}
\setlength{\parindent}{2em}
\setlength{\parskip}{1.25em}
\renewcommand{\baselinestretch}{1.0}
\begin{document}
\begin{center}
{\LARGE \textit{Jones v. Alfred H. Mayer Co.}, 392 U.S. 409 (1968)\par}
\vspace{0.8em}
{\large Case Brief\par}
\end{center}
\vspace{1em}

\section{Memory Jogger}

Bypassing the Fourteenth Amendment's state-action requirement, the Supreme Court held that 42 U.S.C.\ \S~1982---enacted in the Civil Rights Act of 1866---bars purely private racial discrimination in the sale or rental of property and is a valid exercise of Congress's Thirteenth Amendment power to abolish the ``badges and incidents of slavery'' (\href{https://www.courtlistener.com/opinion/107748/jones-v-alfred-h-mayer-co/}{CourtListener}; \href{https://supreme.justia.com/cases/federal/us/392/409/}{Justia}).

\section{Detailed Case Facts}

In June 1965, Joseph Lee Jones, a Black federal employee, and his wife, Barbara Jo Jones, a white federal employee, sought to purchase a lot and have a home built in Paddock Woods, a private suburban housing subdivision in St.\ Louis County, Missouri, being developed by the Alfred H.\ Mayer Co.\ (\href{https://clearinghouse.net/case/10160/}{Civil Rights Litig.\ Clearinghouse}; \href{https://scholarlycommons.law.case.edu/cgi/viewcontent.cgi?params=/context/caselrev/article/2719/&path_info=29_20CaseWResLRev448_1968_1969_.pdf}{Case W.\ Res.\ L.\ Rev.}). Although the couple offered to pay the advertised price, the developer refused to sell them a home solely because Mr.\ Jones was Black (\href{https://clearinghouse.net/case/10160/}{Civil Rights Litig.\ Clearinghouse}).

In September 1965, the Joneses filed suit in federal district court seeking injunctive relief and damages, invoking federal civil rights statutes---principally 42 U.S.C.\ \S~1982, derived from \S~1 of the Civil Rights Act of 1866, 14 Stat.\ 27---which provides that ``[a]ll citizens of the United States shall have the same right, in every State and Territory, as is enjoyed by white citizens thereof to inherit, purchase, lease, sell, hold, and convey real and personal property'' (\href{https://clearinghouse.net/case/10160/}{Civil Rights Litig.\ Clearinghouse}; \href{https://supreme.justia.com/cases/federal/us/392/409/}{Justia}). They also argued in the alternative that the developer's comprehensive control over a planned suburban community amounted to ``state action'' under the Fourteenth Amendment (\href{https://en.wikisource.org/wiki/Jones_v._Mayer/Dissent_Harlan}{Wikisource}). While the case was pending before the Supreme Court in April 1968, Congress enacted the Fair Housing Act (Title VIII of the Civil Rights Act of 1968), 42 U.S.C.\ \S\S~3601--3619, which barred most private housing discrimination prospectively but did not take full effect in time to cover the Joneses' 1965 claim (\href{https://supreme.justia.com/cases/federal/us/392/409/}{Justia}; \href{https://scholarlycommons.law.case.edu/cgi/viewcontent.cgi?params=/context/caselrev/article/2719/&path_info=29_20CaseWResLRev448_1968_1969_.pdf}{Case W.\ Res.\ L.\ Rev.}).

\section{Procedural History}

On June 9, 1966, U.S.\ District Judge John K.\ Regan of the Eastern District of Missouri granted the defendants' motion to dismiss for failure to state a claim, holding that \S~1982 reached only state action and that the developer's private refusal to sell did not implicate the Fourteenth Amendment (\textit{Jones v.\ Alfred H.\ Mayer Co.}, 255 F.\ Supp.\ 115 (E.D.\ Mo.\ 1966)) (\href{https://clearinghouse.net/case/10160/}{Civil Rights Litig.\ Clearinghouse}). On June 26, 1967, the U.S.\ Court of Appeals for the Eighth Circuit, in an opinion by Judge Harry A.\ Blackmun, affirmed the dismissal on the ground that \S~1982 applied only to state action and not to purely private refusals to sell (\textit{Jones v.\ Alfred H.\ Mayer Co.}, 379 F.2d 33 (8th Cir.\ 1967)) (\href{https://clearinghouse.net/case/10160/}{Civil Rights Litig.\ Clearinghouse}; \href{https://supreme.justia.com/cases/federal/us/392/409/}{Justia}). The Supreme Court granted certiorari (389 U.S.\ 968), heard oral argument on April 1--2, 1968, and reversed on June 17, 1968 (\href{https://supreme.justia.com/cases/federal/us/392/409/}{Justia}; \href{https://tile.loc.gov/storage-services/service/ll/usrep/usrep392/usrep392409/usrep392409.pdf}{Library of Congress}).

\section{Judicial Votes}

\textbf{Vote:} 7--2, reversing the judgment of the Eighth Circuit (\href{https://supreme.justia.com/cases/federal/us/392/409/}{Justia}; \href{https://scholarship.law.vanderbilt.edu/cgi/viewcontent.cgi?article=3460&context=vlr}{Ervin, Vand.\ L.\ Rev.}).

\begin{itemize}
\item \textbf{Majority (7):} Justice Potter Stewart (author), joined by Chief Justice Earl Warren and Associate Justices Hugo L.\ Black, William O.\ Douglas, William J.\ Brennan, Jr., Abe Fortas, and Thurgood Marshall.
\item \textbf{Concurrence (1):} Justice William O.\ Douglas filed a separate concurring opinion (while also joining the majority opinion) (\href{https://en.wikisource.org/wiki/Jones_v._Alfred_H._Mayer_Co./Concurrence_Douglas}{Wikisource}).
\item \textbf{Dissent (2):} Justice John Marshall Harlan II (author), joined by Associate Justice Byron R.\ White (\href{https://en.wikisource.org/wiki/Jones_v._Mayer/Dissent_Harlan}{Wikisource}).
\end{itemize}

\section{Holding}

\textbf{Judgment:} Reversed. The Joneses stated a valid federal cause of action under 42 U.S.C.\ \S~1982 against the private developer's racially motivated refusal to sell (\href{https://supreme.justia.com/cases/federal/us/392/409/}{Justia}).

\textbf{Rule Announced:} First, 42 U.S.C.\ \S~1982 prohibits all racial discrimination, private as well as public, in the sale or rental of real or personal property, unaffected by the subsequent enactment of the Fair Housing Act of 1968 (\href{https://supreme.justia.com/cases/federal/us/392/409/}{Justia}). Second, because the Thirteenth Amendment contains no state-action limitation, Section 2 empowers Congress ``rationally to determine what are the badges and the incidents of slavery, and the authority to translate that determination into effective legislation''---including outlawing private racial barriers to acquiring property (\href{https://constitution.findlaw.com/amendment13.html}{FindLaw}). Insofar as \textit{Hodges v.\ United States}, 203 U.S.\ 1 (1906), held otherwise, it was expressly overruled (\href{https://supreme.justia.com/cases/federal/us/392/409/}{Justia}; \href{https://www.law.cornell.edu/constitution-conan/amendment-13/scope-of-enforcement-clause-of-thirteenth-amendment}{Cornell LII}).

\textbf{Holding vs.\ Dicta:} The \emph{holding} establishes (1) \S~1982's reach to purely private racial discrimination in property transactions, (2) Congress's rational-basis authority under Section 2 of the Thirteenth Amendment to regulate private conduct constituting a badge or incident of slavery, and (3) the availability of injunctive relief under \S~1982 (\href{https://supreme.justia.com/cases/federal/us/392/409/}{Justia}). The Court expressly reserved as \emph{non-holdings / dicta} whether the petitioners' alternative Fourteenth Amendment state-action theory succeeded, whether Section 1 of the Thirteenth Amendment \emph{of its own force} (without congressional legislation) prohibits private housing discrimination, and whether compensatory damages (as distinct from punitive damages, which the Court noted were unavailable on the pleaded facts) could be recovered under \S~1982 (\href{https://supreme.justia.com/cases/federal/us/392/409/}{Justia}; \href{https://www.law.cornell.edu/constitution-conan/amendment-13/scope-of-enforcement-clause-of-thirteenth-amendment}{Cornell LII}).

\section{Analysis of Opinions}

\subsection{Majority Opinion (Justice Stewart)}

Justice Stewart began by distinguishing 42 U.S.C.\ \S~1982 from the newly enacted Fair Housing Act of 1968. Whereas Title VIII is a comprehensive open-housing code covering religion and national origin, advertising, financing, and federal administrative enforcement, \S~1982 is a bare, general prohibition limited to racial discrimination in property transactions and enforceable only by private litigation; accordingly, the 1968 Act neither mooted the case nor impliedly narrowed \S~1982 (\href{https://supreme.justia.com/cases/federal/us/392/409/}{Justia}).

Turning to statutory construction, Stewart treated the applicability of \S~1982 to purely private discrimination as an issue of first impression, distinguishing prior statements in \textit{Hurd v.\ Hodge}, 334 U.S.\ 24 (1948), \textit{Corrigan v.\ Buckley}, 271 U.S.\ 323 (1926), and the \textit{Civil Rights Cases}, 109 U.S.\ 3 (1883), as dicta or confined to different questions (\href{https://supreme.justia.com/cases/federal/us/392/409/}{Justia}). On its face, \S~1982 grants all citizens ``the same right'' to purchase and lease property ``as is enjoyed by white citizens,'' a right impaired ``as effectively by `those who place property on the market' as by the State itself'' (\href{http://law2.umkc.edu/faculty/projects/ftrials/conlaw/mayer.html}{UMKC Law}). Examining the 1866 legislative debates, Stewart concluded that the 39th Congress intended \S~1 of the Civil Rights Act of 1866 to protect freedmen not only against state Black Codes but also against private ``custom, or prejudice,'' while deliberately restricting \S~2's criminal penalties to deprivations committed ``under color of law'' (\href{https://opencasebook.org/casebooks/6014-constitutional-law-ii/resources/2.3-jones-v-alfred-h-mayer-co/}{Open Casebooks}). Reenactment of the statute in 1870 after ratification of the Fourteenth Amendment did not shrink its Thirteenth Amendment scope, nor did decades of partial dormancy diminish its force (\href{https://supreme.justia.com/cases/federal/us/392/409/}{Justia}).

On the constitutional question, Stewart invoked the \textit{Civil Rights Cases}, 109 U.S.\ at 20, 23, for the settled principle that the Thirteenth Amendment---unlike the Fourteenth---is an ``absolute declaration'' containing no state-action requirement and empowers Congress to enact ``direct and primary'' legislation operating upon private individuals (\href{https://supreme.justia.com/cases/federal/us/392/409/}{Justia}; \href{https://en.wikipedia.org/wiki/Thirteenth_Amendment_to_the_United_States_Constitution}{Wikipedia}). Under Section 2, Congress possesses rational-basis authority to identify and eliminate the ``badges and the incidents of slavery'' (\href{https://constitution.findlaw.com/amendment13.html}{FindLaw}). Because excluding Black citizens from white communities operated as a private substitute for the Black Codes and ``herds men into ghettos,'' Congress rationally determined that private racial barriers to buying property were a relic of slavery (\href{https://constitution.findlaw.com/amendment13.html}{FindLaw}). Stewart therefore overruled \textit{Hodges v.\ United States} to the extent it conflicted with this principle and did not reach the Joneses' Fourteenth Amendment state-action claim (\href{https://supreme.justia.com/cases/federal/us/392/409/}{Justia}).

\subsection{Concurring Opinion (Justice Douglas)}

Justice Douglas joined the Court's opinion and wrote separately to catalog how the ``badges of slavery'' persisted in American law and private custom a century after emancipation (\href{https://en.wikisource.org/wiki/Jones_v._Alfred_H._Mayer_Co./Concurrence_Douglas}{Wikisource}). Quoting W.E.B.\ Du Bois's \textit{Black Reconstruction in America} and Frederick Douglass's 1881 essay \textit{The Color Line}, Douglas argued that slavery's deepest legacy was the ideology of white supremacy ``in the minds and hearts of many white men'' (\href{https://en.wikisource.org/wiki/Jones_v._Alfred_H._Mayer_Co./Concurrence_Douglas}{Wikisource}). Surveying Supreme Court decisions involving voting, juries, schools, interracial marriage, public accommodations, and housing, Douglas concluded that residential segregation was a pervasive ``custom'' of slavery that the 1866 Congress rightly sought to eradicate under \S~1982 (\href{https://en.wikisource.org/wiki/Jones_v._Alfred_H._Mayer_Co./Concurrence_Douglas}{Wikisource}).

\subsection{Dissenting Opinion (Justice Harlan, joined by Justice White)}

Justice Harlan condemned the majority's ruling as ``most ill-considered and ill-advised'' (\href{https://supreme.justia.com/cases/federal/us/392/409/}{Justia}; \href{https://en.wikisource.org/wiki/Jones_v._Mayer/Dissent_Harlan}{Wikisource}). First, Harlan argued that the Court's statutory construction was ``almost surely wrong'' (\href{http://law2.umkc.edu/faculty/projects/ftrials/conlaw/mayer.html}{UMKC Law}). In his reading of the 1866 debates, Congress's phrase ``the same right'' meant equal \emph{legal capacity} to buy and own property free from discriminatory state laws and official customs (such as the Black Codes), not a substantive right to compel an unwilling private owner to sell; moreover, the 1866 text originally qualified \S~1 with ``any law, statute, ordinance, regulation, or custom, to the contrary notwithstanding,'' and \S~2's ``color of law'' language confirmed that both sections targeted state action (\href{https://en.wikisource.org/wiki/Jones_v._Mayer/Dissent_Harlan}{Wikisource}; \href{https://scholarship.law.vanderbilt.edu/cgi/viewcontent.cgi?article=3460&context=vlr}{Ervin, Vand.\ L.\ Rev.}).

Second, Harlan emphasized that prior Supreme Court decisions---including the \textit{Civil Rights Cases}, \textit{Virginia v.\ Rives}, \textit{Corrigan v.\ Buckley}, and \textit{Hurd v.\ Hodge}---had consistently understood \S~1982 to reach only governmental action, and even the 1866 Act's supporters had doubted Congress's power under the Thirteenth Amendment alone, prompting them to pass the Fourteenth Amendment (\href{http://law2.umkc.edu/faculty/projects/ftrials/conlaw/mayer.html}{UMKC Law}; \href{https://en.wikisource.org/wiki/Jones_v._Mayer/Dissent_Harlan}{Wikisource}). Finally, because Congress had just enacted the carefully calibrated Fair Housing Act of 1968---which included explicit exemptions such as small owner-occupied dwellings---Harlan urged that the Court should have dismissed the writ of certiorari as improvidently granted rather than judicially fashioning a broader, uncalibrated fair-housing mandate out of a century-old statute (\href{https://en.wikisource.org/wiki/Jones_v._Mayer/Dissent_Harlan}{Wikisource}; \href{https://scholarlycommons.law.case.edu/cgi/viewcontent.cgi?params=/context/caselrev/article/2719/&path_info=29_20CaseWResLRev448_1968_1969_.pdf}{Case W.\ Res.\ L.\ Rev.}).

\section{Examples: Future Applications}

\begin{enumerate}
\item \textbf{Same-Side (Covered / Valid): Private Commercial School Excluding Black Applicants.} A privately operated, nonsectarian day school that advertises to the general public refuses to admit Black students solely because of their race. Under \textit{Jones}'s reading of the Civil Rights Act of 1866 and Congress's Thirteenth Amendment power, 42 U.S.C.\ \S~1981 (the companion contract-rights provision from \S~1 of the 1866 Act) validly bars such private racial refusals to contract without requiring state action (as the Court held in \textit{Runyon v.\ McCrary}, 427 U.S.\ 160 (1976)) (\href{https://law.justia.com/constitution/us/amendment-13/01-abolition-of-slavery.html}{Justia Constitution Annotated}).

\item \textbf{Same-Side (Covered / Valid): Federal Statute Criminalizing Racially Motivated Private Hate Crimes.} Congress enacts a statute making it a federal crime for private individuals to commit racially motivated violent assaults, even without any nexus to interstate commerce or state officials. Under \textit{Jones}, Congress may rationally determine that private racially motivated violence is a badge or incident of slavery and outlaw it directly under Section 2 of the Thirteenth Amendment (as federal courts upheld for 18 U.S.C.\ \S~249(a)(1) in \textit{United States v.\ Hatch}, 722 F.3d 1193 (10th Cir.\ 2013), and \textit{United States v.\ Nelson}, 277 F.3d 164 (2d Cir.\ 2002)) (\href{https://caselaw.findlaw.com/court/us-10th-circuit/1637502.html}{FindLaw}; \href{https://caselaw.findlaw.com/court/us-2nd-circuit/1333018.html}{FindLaw}).

\item \textbf{Opposite-Side (Not Covered): Private Religious or Sex Discrimination Under \S~1982.} A private homeowner refuses to rent a basement apartment to a prospective tenant solely because of the tenant's religion or sex, with no racial or ancestral animus. Because \textit{Jones} expressly held that \S~1982 ``deals only with racial discrimination and does not address itself to discrimination on grounds of religion or national origin'' (or sex), \S~1982 provides no remedy (though the Fair Housing Act may apply if no statutory exemption covers the landlord) (\href{https://supreme.justia.com/cases/federal/us/392/409/}{Justia}; \href{https://dc.law.utah.edu/cgi/viewcontent.cgi?article=2245&context=ulr}{Utah L.\ Rev.}).

\item \textbf{Opposite-Side (Not Covered): Direct Judicial Enforcement of Section 1 Without Congressional Legislation.} Residents sue a municipality directly under Section 1 of the Thirteenth Amendment---without invoking any federal statute---challenging a facially neutral street closure that reduces traffic through a white neighborhood while inconveniencing Black commuters. In \textit{City of Memphis v.\ Greene}, 451 U.S.\ 100 (1981), the Court declined to treat such a neutral municipal measure as a judicially enforceable ``badge of slavery'' under Section 1 alone, underscoring that \textit{Jones}'s broad ``badges and incidents'' power belongs primarily to \emph{Congress} under Section 2 (\href{https://www.law.cornell.edu/constitution-conan/amendment-13/scope-of-enforcement-clause-of-thirteenth-amendment}{Cornell LII}).

\item \textbf{Fence-Sitter (Close Case): Federal Thirteenth Amendment Legislation Targeting Non-Racial Private Discrimination.} Relying exclusively on Section 2 of the Thirteenth Amendment (and not the Commerce Clause), Congress passes a statute banning private landlords from discriminating based on sexual orientation or disability. Proponents would cite \textit{Jones}'s deferential rational-basis test---asking whether Congress could rationally deem status-based subordination a modern ``badge'' of involuntary servitude (\href{https://scholars.law.unlv.edu/cgi/viewcontent.cgi?article=1013&context=amend13_bi}{UNLV Scholarly Commons}). Opponents would counter that \textit{Jones} anchored the ``badges and incidents'' doctrine to the historical institution of chattel slavery and racial caste, and that extending Section 2 to non-racial private conduct without a concrete historical link to slavery risks turning the Thirteenth Amendment into an unbounded federal police power---a tension courts have flagged in light of \textit{City of Boerne v.\ Flores} (\href{https://caselaw.findlaw.com/court/us-10th-circuit/1637502.html}{FindLaw}; \href{https://lawreview.law.ucdavis.edu/sites/g/files/dgvnsk15026/files/media/documents/DavisVol40No4_Carter.pdf}{Carter, U.C.\ Davis L.\ Rev.}).
\end{enumerate}

\section{Critique}

\textbf{Originalist, Textualist, and Historical Perspectives:} Originalist and historical critics---echoing Justice Harlan's dissent---have long faulted \textit{Jones} for misreading the 1866 legislative record. In a famous contemporary critique, Gerhard Casper argued that the Court treated history as a ``bemused and confused muse,'' while Senator Sam Ervin attacked the decision as ``judicial activism run riot'' that severed \S~1 from \S~2's ``color of law'' restriction (\href{https://caselaw.findlaw.com/court/us-10th-circuit/1637502.html}{FindLaw}; \href{https://scholarship.law.vanderbilt.edu/cgi/viewcontent.cgi?article=3460&context=vlr}{Ervin, Vand.\ L.\ Rev.}). Modern textualists and scholars such as Jennifer Mason McAward and George Rutherglen also warn that \textit{Jones}'s deferential rational-basis standard under Section 2 sits in uneasy tension with the Supreme Court's later ``congruence and proportionality'' limit on Fourteenth Amendment enforcement power in \textit{City of Boerne v.\ Flores}, 521 U.S.\ 507 (1997), risking an unchecked federal police power unless ``badges and incidents'' is tethered to historical chattel slavery (\href{https://scholars.law.unlv.edu/cgi/viewcontent.cgi?article=1013&context=amend13_bi}{UNLV Scholarly Commons}; \href{https://caselaw.findlaw.com/court/us-10th-circuit/1637502.html}{FindLaw}).

\textbf{Progressive and Doctrinal Perspectives:} Progressive and Reconstruction-history scholars strongly defend both \textit{Jones}'s result and its constitutional foundation. Darrell A.\ H.\ Miller demonstrates that the 39th Congress understood private ``white cartels'' and customary boycotts to be central mechanisms recreating slavery after emancipation, vindicating Stewart's reading of the 1866 Act (\href{https://scholarship.law.duke.edu/cgi/viewcontent.cgi?article=5768&context=faculty_scholarship}{Miller, Fordham L.\ Rev.}). Scholars such as William M.\ Carter, Jr., and Alexander Tsesis celebrate \textit{Jones} for reviving the Thirteenth Amendment as an affirmative charter of substantive freedom that empowers Congress to dismantle systemic private racial subordination without distorting the Commerce Clause or confronting the Fourteenth Amendment's state-action barrier (\href{https://lawreview.law.ucdavis.edu/sites/g/files/dgvnsk15026/files/media/documents/DavisVol40No4_Carter.pdf}{Carter, U.C.\ Davis L.\ Rev.}; \href{https://scholars.law.unlv.edu/cgi/viewcontent.cgi?article=1013&context=amend13_bi}{UNLV Scholarly Commons}).

\textbf{Underlying Competing Values:} \textit{Jones} sits at the crossroads of four competing constitutional values: (1) eradicating private racial caste and residential segregation versus preserving private property autonomy and freedom of association, and (2) empowering Congress to fulfill Reconstruction's promise of universal freedom versus maintaining federalism limits and judicial restraint when Congress has already crafted a narrower modern statutory scheme (\href{https://www.law.cornell.edu/constitution-conan/amendment-13/scope-of-enforcement-clause-of-thirteenth-amendment}{Cornell LII}; \href{https://www.scotusblog.com/2010/02/jones-v-alfred-mayer-and-the-uniqueness-of-race/}{SCOTUSblog}).

\section{Key Quotations}

\begin{quote}
``We hold that \S~1982 bars all racial discrimination, private as well as public, in the sale or rental of property, and that the statute, thus construed, is a valid exercise of the power of Congress to enforce the Thirteenth Amendment.''
\end{quote}
\noindent --- Justice Stewart, \textit{Jones v.\ Alfred H.\ Mayer Co.}, 392 U.S.\ 409, 413 (1968) (\href{https://supreme.justia.com/cases/federal/us/392/409/}{Justia}).

\begin{quote}
``Surely Congress has the power under the Thirteenth Amendment rationally to determine what are the badges and the incidents of slavery, and the authority to translate that determination into effective legislation.''
\end{quote}
\noindent --- Justice Stewart, \textit{Jones v.\ Alfred H.\ Mayer Co.}, 392 U.S.\ 409, 440 (1968) (\href{https://constitution.findlaw.com/amendment13.html}{FindLaw}).

\begin{quote}
``At the very least, the freedom that Congress is empowered to secure under the Thirteenth Amendment includes the freedom to buy whatever a white man can buy, the right to live wherever a white man can live. If Congress cannot say that being a free man means at least this much, then the Thirteenth Amendment made a promise the Nation cannot keep.''
\end{quote}
\noindent --- Justice Stewart, \textit{Jones v.\ Alfred H.\ Mayer Co.}, 392 U.S.\ 409, 443 (1968) (\href{https://supreme.justia.com/cases/federal/us/392/409/}{Justia}).

\begin{quote}
``Some badges of slavery remain today. While the institution has been outlawed, it has remained in the minds and hearts of many white men.''
\end{quote}
\noindent --- Justice Douglas, concurring, \textit{Jones v.\ Alfred H.\ Mayer Co.}, 392 U.S.\ 409, 445 (1968) (\href{https://en.wikisource.org/wiki/Jones_v._Alfred_H._Mayer_Co./Concurrence_Douglas}{Wikisource}).

\begin{quote}
``The decision in this case appears to me to be most ill-considered and ill-advised. \dots\ For reasons which follow, I believe that the Court's construction of \S~1982 as applying to purely private action is almost surely wrong, and at the least is open to serious doubt.''
\end{quote}
\noindent --- Justice Harlan, dissenting, \textit{Jones v.\ Alfred H.\ Mayer Co.}, 392 U.S.\ 409, 449--50 (1968) (\href{https://supreme.justia.com/cases/federal/us/392/409/}{Justia}; \href{http://law2.umkc.edu/faculty/projects/ftrials/conlaw/mayer.html}{UMKC Law}).


\end{document}
